\chapter{Van eicel tot geboorte}\label{ch:g8:from-egg-to-birth}

\Cref{ch:g5:human-reproduction} vertelde de negen maanden als een verhaal;
inmiddels heb je elk instrument dat nodig is om dat verhaal goed te
spelen: cellen en hun delingen, de twee \omterm{def:g8:sexual-reproduction:gametes}{gameten} en hun ontmoetingsplaats,
\omterm{def:g8:puberty:hormones}{hormonen}, de leveringen van het \omterm{prop:g4:journey-of-food:absorb}{bloed}. Dit hoofdstuk speelt de menselijke
ontwikkeling opnieuw af met de machinerie in beeld --- van één \omterm{def:g5:first-look-at-cells:cell}{cel} tot
drieduizend miljard, in veertig weken.

\section{De eerste week: van zygote tot huurder}

\begin{proposition}[De reis door de eileider omlaag]\label{prop:g8:from-egg-to-birth:firstweek}
De \omterm{def:g5:flower-to-fruit:steps}{bevruchting} in de \omterm{def:g8:reproductive-systems:female}{eileider}
(\cref{def:g8:reproductive-systems:female}) start twee klokken tegelijk:
\begin{enumerate}
  \item de \emph{deling}: de \omterm{prop:g8:sexual-reproduction:core}{zygote} deelt zich binnen een dag, en dan
        keer op keer (\cref{ex:g5:first-look-at-cells:one}) --- twee,
        vier, acht cellen, en tegen dag vier een moerbeibolletje van
        tientallen, en dat allemaal terwijl ze door de \omterm{def:g8:reproductive-systems:female}{eileider} naar de
        \omterm{prop:g5:human-reproduction:plan}{baarmoeder} afdrijft;
  \item de \emph{innesteling}: rond dag zeven bereikt het bolletje de
        \omterm{prop:g5:human-reproduction:plan}{baarmoeder} en nestelt het zich in het klaargemaakte slijmvlies
        (gebeurtenis 4 uit
        \cref{prop:g8:reproductive-systems:cycle}) --- de huurder komt
        aan, de cyclus stopt, en \omterm{def:g8:puberty:hormones}{hormonen} seinen het hele stelsel in:
        \omterm{def:g5:human-reproduction:pregnancy}{zwangerschap}.
\end{enumerate}
\end{proposition}
\admitted

\begin{example}[Waar de eerste delingen van leven]\label{ex:g8:from-egg-to-birth:provisions}
De hele eerste week eet de reiziger geen maaltijd en krijgt hij geen
levering: hij draait op de eigen proviand van de \omterm{def:g4:animal-reproduction:sexual}{eicel} --- precies
waarvoor de voorraad uit \cref{ex:g8:reproductive-systems:egg} bedoeld
was. De \omterm{def:g8:reproductive-systems:male}{zaadcel} bracht de \omterm{def:g6:cell-unit-of-life:parts}{celkern} van de vader mee en verder weinig; de
\omterm{def:g4:animal-reproduction:sexual}{eicel} bracht de andere \omterm{def:g6:cell-unit-of-life:parts}{celkern} \emph{én} de voorraadkast. De innesteling
is, onder meer, de voorraadkast die opraakt en de huurder die zich op de
roomservice aansluit.
\end{example}

\section{Embryo, placenta, foetus}

\begin{definition}[De placenta]\label{def:g8:from-egg-to-birth:placenta}
Uit het ingenestelde bolletje groeien tegelijk twee bouwwerken: het
\emph{\omterm{def:g5:human-reproduction:pregnancy}{embryo}} zelf, en zijn levensondersteunende koppeling --- de
\emph{placenta}\index{placenta}, een schijf die in de wand van de
\omterm{prop:g5:human-reproduction:plan}{baarmoeder} groeit en waar het \omterm{prop:g4:journey-of-food:absorb}{bloed} van de moeder en dat van het \omterm{def:g5:human-reproduction:pregnancy}{embryo}
innig dicht bij elkaar komen \emph{zonder zich te mengen}. Over haar
uitgestrekte dunne oppervlakken --- opnieuw het bestek uit
\cref{prop:g7:how-animals-breathe:spec} --- steken \omterm{def:g7:digestion-and-nutrients:families}{voedingsstoffen} en
\omterm{def:g7:what-is-respiration:respiration}{zuurstof} over naar het \omterm{def:g5:human-reproduction:pregnancy}{embryo}, en afvalstoffen en \omterm{def:g7:what-is-respiration:respiration}{koolstofdioxide} terug
(de \omterm{prop:g5:human-reproduction:cord}{navelstreng} uit \cref{prop:g5:human-reproduction:cord} loopt tussen
placenta en \omterm{def:g5:human-reproduction:pregnancy}{embryo}). De placenta is \omterm{def:g7:how-animals-breathe:four}{long}, darm en \omterm{def:g7:removing-wastes:kidneys}{nier} voor de passagier
--- alle drie de grenzen in één orgaan.
\end{definition}

\begin{proposition}[Het tijdschema, met instrumenten herzien]\label{prop:g8:from-egg-to-birth:timetable}
De mijlpalen uit \cref{ex:g5:human-reproduction:timetable}, nu verklaard:
\begin{itemize}
  \item week 1--8, het \emph{\omterm{def:g5:human-reproduction:pregnancy}{embryo}}: het bouwplan wordt uitgezet ---
        cellen delen zich, verhuizen en \emph{specialiseren} zich
        (\cref{rem:g6:cell-unit-of-life:twoways}) tot \omterm{def:g4:heart-and-blood:heart}{hart}, \omterm{def:g5:brain-in-command:system}{zenuw}, \omterm{def:g3:bones-and-muscles:skeleton}{bot} en
        de rest; het \omterm{def:g4:heart-and-blood:heart}{hart} klopt vanaf week vier, omdat de huishouding van
        voorraadkast en \omterm{def:g8:from-egg-to-birth:placenta}{placenta} haar pomp vroeg nodig heeft;
  \item week 9--40, de \emph{\omterm{def:g5:human-reproduction:pregnancy}{foetus}}: het plan gebouwd, en daarna groeien
        en rijpen --- de \omterm{def:g5:brain-in-command:system}{hersenen} die zichzelf bedraden, de \omterm{def:g7:how-animals-breathe:four}{longen} die leeg
        oefenen (\cref{exo:g5:human-reproduction:11} verklaarde hun debuut
        bij de eerste kreet), en de schopjes van maand vier die werkende
        \omterm{def:g3:bones-and-muscles:muscle}{spieren} aankondigen;
  \item de hele tijd door: \omterm{def:g8:puberty:hormones}{hormonen} houden de \omterm{prop:g5:human-reproduction:plan}{baarmoeder} stil en de cyclus
        opgeschort --- de scheikunde uit
        \cref{def:g8:puberty:hormones}, die de administratie van de
        \omterm{def:g5:human-reproduction:pregnancy}{zwangerschap} voert.
\end{itemize}
\end{proposition}
\admitted

\begin{example}[Waarom de vroege weken de meeste zorg vragen]\label{ex:g8:from-egg-to-birth:earlycare}
\Cref{prop:g5:human-reproduction:care} waarschuwde dat tabak en alcohol
naar de \omterm{def:g3:stages-of-life:stages}{baby} oversteken; het tijdschema zegt wanneer dat het meest telt:
in de embryoweken waarin het plan wordt uitgezet, waarin elk orgaan een
klein groepje specialiserende cellen is en één verstoorde instructie door
de hele bouw kan doorklinken --- en dat zijn precies de weken waarin een
\omterm{def:g5:human-reproduction:pregnancy}{zwangerschap} onopgemerkt kan blijven. Vandaar de botte rekensom van de
artsen: de zorg begint niet bij de eerste echo maar bij de eerste
mogelijkheid.
\end{example}

\section{De geboorte, en de overdracht}

\begin{proposition}[De geboorte als een door hormonen bevolen overdracht]\label{prop:g8:from-egg-to-birth:birth}
Rond week veertig beëindigen hormoonsignalen --- die van de \omterm{def:g3:stages-of-life:stages}{baby} en die
van de moeder in gesprek --- de lange stilte van de \omterm{prop:g5:human-reproduction:plan}{baarmoeder}:
\begin{enumerate}
  \item de \emph{bevalling}: de \omterm{prop:g5:human-reproduction:plan}{baarmoeder}, de sterkste \omterm{def:g3:bones-and-muscles:muscle}{spier} van het
        lichaam (\cref{rem:g8:reproductive-systems:living} maakte kennis
        met haar zachtere werk), trekt in steeds sterkere golven samen die
        de weg openen en de \omterm{def:g3:stages-of-life:stages}{baby} door het geboortekanaal omlaag duwen;
  \item de \emph{\omterm{def:g2:born-grow-die:cycle}{geboorte}}: meestal met het \omterm{def:g1:my-body:parts}{hoofd} eerst
        (\cref{ex:g5:human-reproduction:birth});
  \item de \emph{overdracht}: de eerste kreet vult de nooit gebruikte
        \omterm{def:g7:how-animals-breathe:four}{longen}; het verkeer door de streng stopt en de streng wordt
        doorgeknipt; binnen enkele minuten draait de pasgeborene op zijn
        eigen \omterm{def:g7:how-animals-breathe:four}{longen}, op zijn darm-in-wording en op zijn \omterm{def:g7:removing-wastes:kidneys}{nieren} --- op
        elke grens die de \omterm{def:g8:from-egg-to-birth:placenta}{placenta} gedekt had
        (\cref{exo:g5:human-reproduction:11}, nu met de volledige kaart);
  \item de \emph{nageboorte}: de \omterm{def:g8:from-egg-to-birth:placenta}{placenta}, haar werk gedaan, laat los en
        volgt --- het enige orgaan dat het lichaam voor één keer gebruik
        bouwt en daarna weggooit.
\end{enumerate}
\end{proposition}
\admitted

\begin{example}[De eerste doorlichting van de pasgeborene]\label{ex:g8:from-egg-to-birth:audit}
Verloskundigen geven een pasgeborene een paar minuten na de \omterm{def:g2:born-grow-die:cycle}{geboorte} een
score --- \omterm{def:g7:what-is-respiration:respiration}{ademhaling}, \omterm{def:g4:heart-and-blood:heart}{hartslag}, kleur, beweging, huilen. Lees die
checklist met de \omterm{def:g1:the-five-senses:sense}{ogen} van dit boek: het is de lus van aanvoer en afvoer
uit \cref{prop:g7:removing-wastes:complete}, doorgelicht bij haar eerste
solovlucht. Het examen dat elk mens als eerste aflegt, is een
fysiologische test.
\end{example}

\begin{method}[Elk stadium van de zwangerschap uitleggen]\label{met:g8:from-egg-to-birth:explain}
Antwoord op elke vraag over een moment uit de negen maanden met vier
instrumenten:
\begin{enumerate}
  \item \emph{cellen}: wat doen de delingen en specialisaties?
  \item \emph{aanvoer}: voorraadkast of \omterm{def:g8:from-egg-to-birth:placenta}{placenta} --- wat steekt er over,
        en wat mag juist niet?
  \item \emph{\omterm{def:g8:puberty:hormones}{hormonen}}: wat wordt er vastgehouden, of bevolen?
  \item \emph{zorg}: wat vraagt dit stadium van de gewoonten van de moeder
        en van de kalender van de artsen?
\end{enumerate}
\end{method}

\section{Oefeningen}

\begin{exercise}[$\star$]\label{exo:g8:from-egg-to-birth:1}
Zet de eerste week op de klok: het begin van de deling, het
moerbeibolletje, de innesteling.
\end{exercise}

\begin{exercise}[$\star$]\label{exo:g8:from-egg-to-birth:2}
Waar leeft de reiziger van de eerste week van, en wat maakt een eind aan
die regeling?
\end{exercise}

\begin{exercise}[$\star$]\label{exo:g8:from-egg-to-birth:3}
Wat is de \omterm{def:g8:from-egg-to-birth:placenta}{placenta} --- wat steekt er in elke richting over, en wat steekt
er nooit over?
\end{exercise}

\begin{exercise}[$\star$]\label{exo:g8:from-egg-to-birth:4}
Deel de veertig weken in \omterm{def:g5:human-reproduction:pregnancy}{embryo} en \omterm{def:g5:human-reproduction:pregnancy}{foetus} in, met de bezigheid van elke
fase.
\end{exercise}

\begin{exercise}[$\star$]\label{exo:g8:from-egg-to-birth:5}
Waarom klopt het \omterm{def:g4:heart-and-blood:heart}{hart} vanaf week vier, als eerste van alle organen?
\end{exercise}

\begin{exercise}[$\star$]\label{exo:g8:from-egg-to-birth:6}
Noem de vier genummerde gebeurtenissen van de \omterm{def:g2:born-grow-die:cycle}{geboorte} op volgorde.
\end{exercise}

\begin{exercise}[$\star\star$]\label{exo:g8:from-egg-to-birth:7}
Welke taken van drie organen neemt de \omterm{def:g8:from-egg-to-birth:placenta}{placenta} waar? Noem bij elk het
grenshoofdstuk dat erover ging.
\end{exercise}

\begin{exercise}[$\star\star$]\label{exo:g8:from-egg-to-birth:8}
Waarom vragen de embryoweken de meeste zorg --- twee redenen uit
\cref{ex:g8:from-egg-to-birth:earlycare}?
\end{exercise}

\begin{exercise}[$\star\star$]\label{exo:g8:from-egg-to-birth:9}
Lees het scoreformulier van de pasgeborene als de lus uit
\cref{prop:g7:removing-wastes:complete}: koppel \omterm{def:g7:what-is-respiration:respiration}{ademhaling}, \omterm{def:g4:heart-and-blood:heart}{hartslag} en
kleur aan hun stelsels.
\end{exercise}

\begin{exercise}[$\star\star$]\label{exo:g8:from-egg-to-birth:10}
De eeneiige tweeling opnieuw bekeken
(\cref{rem:g5:human-reproduction:twins}): op welk moment van
\cref{prop:g8:from-egg-to-birth:firstweek} gebeurt de splitsing van de ene
\omterm{prop:g8:sexual-reproduction:core}{zygote}, en waarom maakt die timing hen ``gelijk als kopieën''?
\end{exercise}

\begin{exercise}[$\star\star$]\label{exo:g8:from-egg-to-birth:11}
Laat \cref{met:g8:from-egg-to-birth:explain} lopen over maand vier --- de
maand van de eerste schopjes.
\end{exercise}

\begin{exercise}[$\star\star\star$]\label{exo:g8:from-egg-to-birth:12}
``De \omterm{def:g2:born-grow-die:cycle}{geboorte} is een wisseling van grenzen.'' Onderbouw dat: noem elke
uitwisseling die de \omterm{def:g8:from-egg-to-birth:placenta}{placenta} afhandelde, en het orgaan dat elk daarvan bij
de eerste kreet overneemt --- met de timing op minuutschaal die de
overdracht zo opmerkelijk maakt.
\end{exercise}

\section{Probleem: De tijdlijn van de verloskundige}

\begin{problem}[{Weekendprobleem --- een zwangerschap gevolgd, bezoek voor
bezoek}]\label{pb:g8:from-egg-to-birth:1}
Het dossier van een verloskundige voor één (verzonnen) \omterm{def:g5:human-reproduction:pregnancy}{zwangerschap}: een
intakebezoek in week 8, echo's in week 12 en 20, controles in week 28 en
36, en de \omterm{def:g2:born-grow-die:cycle}{geboorte} in week 40. Lees het dossier met de instrumenten van
dit hoofdstuk.

\medskip\noindent\textbf{Deel I --- De eerste bladzijden van het
dossier.}
\begin{enumerate}
  \item De \omterm{def:g5:human-reproduction:pregnancy}{zwangerschap} werd weken vóór het eerste bezoek met een
        hormoontest bevestigd. Welk signaal leest zo'n test, en vanaf welke
        gebeurtenis uit
        \cref{prop:g8:from-egg-to-birth:firstweek} bestaat het?
  \item In week 8 noteert de verloskundige ``\omterm{def:g5:human-reproduction:pregnancy}{embryo} compleet in plan,
        ongeveer zo groot als een boon''. Vertaal dat in de termen van het
        tijdschema --- welke fase sluit hier af?
  \item Het advies bij het intakebezoek begint met geen alcohol, geen rook
        --- en foliumzuur, een bouwvitamine, ``het liefst al vóór het
        begin''. Onderbouw alle drie met
        \cref{ex:g8:from-egg-to-birth:earlycare}.
  \item De echo van week 12 laat de \omterm{def:g4:heart-and-blood:heart}{hartslag} zien die in week 4 begon.
        Waarom zo vroeg, in termen van aanvoer?
\end{enumerate}

\medskip\noindent\textbf{Deel II --- De middelste maanden.}
\begin{enumerate}[resume]
  \item De echo van week 20 meet elk orgaan en de ligging van de \omterm{def:g8:from-egg-to-birth:placenta}{placenta}.
        Geef de bezigheid van de foetusfase, en de drie taken van de
        \omterm{def:g8:from-egg-to-birth:placenta}{placenta} die gecontroleerd worden.
  \item De moeder meldt de eerste schopjes rond week 18. Welke machinerie
        kondigt zich aan (\cref{prop:g5:brain-in-command:loop} noemt de
        keten)?
  \item In week 28 controleert de verloskundige het \omterm{prop:g4:journey-of-food:absorb}{bloed} van de moeder op
        suiker en druk. Waarom is haar bloedsomloop dubbel belast ---
        wiens grenzen draait ze?
  \item Week 30--36: de \omterm{def:g5:human-reproduction:pregnancy}{foetus} groeit vooral, en de \omterm{def:g7:how-animals-breathe:four}{longen} oefenen.
        Oefenen waarvoor --- en voor welk enkel moment?
\end{enumerate}

\medskip\noindent\textbf{Deel III --- De laatste bladzijde.}
\begin{enumerate}[resume]
  \item Week 40: de bevalling. Noem de \omterm{def:g3:bones-and-muscles:muscle}{spier}, haar twee taken, en wie het
        werk bevolen heeft
        (\cref{prop:g8:from-egg-to-birth:birth}).
  \item De geboorteregel in het dossier: ``meteen krachtig gehuil; streng
        afgeklemd en doorgeknipt; \omterm{def:g8:from-egg-to-birth:placenta}{placenta} compleet''. Licht de drie
        zinsdelen door met de overdracht.
  \item De pasgeborene scoort hoog op de checklist van de eerste minuten.
        Zeg wat die checklist bevestigt, in de termen van
        \cref{ex:g8:from-egg-to-birth:audit}.
  \item Sluit het dossier af met één zin voor de ouders: wat de veertig
        weken gebouwd hebben, uit welke ene \omterm{def:g5:first-look-at-cells:cell}{cel}, en met welke machinerie.
\end{enumerate}
\end{problem}
